% Toolbox Guide of the Self-Study: Learn and Practise Toolbox
% Compile with: julia -e 'using PollisTeX; display(compile("guide.tex"))'
% The look of a Pollis Toolbox Guide: that of the Toolbox Vignettes (vignette_style.tex), with its own name.
% Set \vignetteauthor, \vignettedate and, if not CC BY-SA 4.0, \vignettelicense after % The look of a Pollis Toolbox Guide: that of the Toolbox Vignettes (vignette_style.tex), with its own name.
% Set \vignetteauthor, \vignettedate and, if not CC BY-SA 4.0, \vignettelicense after % The look of a Pollis Toolbox Guide: that of the Toolbox Vignettes (vignette_style.tex), with its own name.
% Set \vignetteauthor, \vignettedate and, if not CC BY-SA 4.0, \vignettelicense after \input{guide_style}.
\documentclass[10pt,a4paper]{article}

\usepackage[a4paper,left=18mm,right=18mm,top=15mm,bottom=17mm,footskip=8mm]{geometry}
\usepackage{iftex}
\ifPDFTeX\usepackage[T1]{fontenc}\fi
\usepackage[default,semibold]{sourcesanspro}
\usepackage{textcomp}
\usepackage[italic]{mathastext}
\usepackage{amsmath}
\usepackage{microtype}
\usepackage[table]{xcolor}
\usepackage{tikz}
\usepackage[most]{tcolorbox}
\usepackage{fontawesome5}
\usepackage{tabularx,booktabs}
\usepackage{enumitem}
\usepackage{fancyhdr}
\usepackage[hidelinks]{hyperref}

% ---------------------------------------------------------------- colours (as in the brochure)
\definecolor{navy}{HTML}{1C3A5E}
\definecolor{pblue}{HTML}{3B78B5}
\definecolor{amber}{HTML}{F2B33D}
\definecolor{deepamber}{HTML}{D9902A}
\definecolor{pale}{HTML}{EEF3F9}
\definecolor{ink}{HTML}{2A2F36}
\definecolor{logonavy}{HTML}{173D63}
\definecolor{logodark}{HTML}{C98A1E}
\definecolor{logoamber}{HTML}{F2B446}
\definecolor{logocream}{HTML}{FEEFD2}
\definecolor{soft}{HTML}{5E6773}
\definecolor{agree}{HTML}{2F7D4F}
\definecolor{differ}{HTML}{B23B3B}
\color{ink}

\setlength{\parindent}{0pt}
\setlength{\parskip}{4pt}
\setlist{nosep,leftmargin=*,itemsep=2pt}
\renewcommand{\arraystretch}{1.15}

\pagestyle{fancy}
\fancyhf{}
\renewcommand{\headrulewidth}{0pt}
\fancyfoot[L]{\footnotesize\color{soft}\textbf{Pollis}\enspace Toolbox Guide\enspace·\enspace\nmtheme}
\fancyfoot[R]{\footnotesize\color{soft}\thepage}
\newcommand{\nmtheme}{}
\newcommand{\vignetteauthor}{}
\newcommand{\vignettedate}{\today}
\newcommand{\vignettelicense}{\href{https://creativecommons.org/licenses/by-sa/4.0/}{CC BY-SA 4.0}}

% The Pollis logo, redrawn so it prints sharply at any size
\newcommand{\pollislogo}[1]{%
  \begin{tikzpicture}[x=#1/512,y=-#1/512,baseline=(current bounding box.center)]
    \fill[logonavy] (0,0) rectangle (512,512);
    \fill[logodark] (36,422) -- (104,256) -- (172,422) -- cycle;
    \fill[logoamber] (104,422) -- (203,180) -- (302,422) -- cycle;
    \fill[logocream] (203,422) -- (339,89) -- (476,422) -- cycle;
  \end{tikzpicture}}

% Top of the first page: logo, name; then the kicker, title and a one-line promise
\newcommand{\nmhead}[3]{%
  \renewcommand{\nmtheme}{#1}%
  \thispagestyle{fancy}%
  \begin{minipage}[c]{13mm}\pollislogo{11mm}\end{minipage}%
  \begin{minipage}[c]{\dimexpr\linewidth-13mm\relax}
    {\color{navy}\Large\bfseries Pollis}\enspace{\color{soft}Toolbox Guide}\hfill
    {\color{deepamber}\small\bfseries\MakeUppercase{#1}}
  \end{minipage}\par\vspace{7mm}
  {\color{navy}\fontsize{22}{26}\selectfont\bfseries #2}\par\vspace{4pt}
  {\color{soft}\large #3}\par\vspace{5mm}}

\newcommand{\nmsection}[1]{\par\vspace{5pt}{\color{navy}\large\bfseries #1}\par\vspace{1pt}}
\newcommand{\nmsub}[1]{\par\vspace{2pt}{\color{deepamber}\bfseries #1}\par}
\newcommand{\gt}{\,{\color{deepamber}\bfseries ›}\,}
\newcommand{\menupath}[1]{{\bfseries\color{navy}#1}}
\newcommand{\fn}[1]{\texttt{\small #1}}
% Citations: \nmcite{key}{Breiman (2001)} links to the entry \nmref{key}{...} in the References
\newcommand{\nmcite}[2]{\hyperlink{ref:#1}{\color{pblue}#2}}
\newcommand{\nmref}[2]{\raisebox{1.1em}[0pt][0pt]{\hypertarget{ref:#1}{}}#2\par}
\newcommand{\nmdoi}[1]{\href{https://doi.org/#1}{\color{pblue}doi:#1}}
% Where the package is: its name and a link (its GitHub repository, once published)
\newcommand{\package}[2]{\href{#2}{\fn{#1.jl}}}

% Key facts at the top: source, missing in Julia, package, result
\newcommand{\factsheet}[4]{%
  \begin{tcolorbox}[enhanced, colback=pale, colframe=pale, boxrule=0pt, sharp corners,
    left=7pt, right=7pt, top=5pt, bottom=5pt, fontupper=\small]
  \begin{tabularx}{\linewidth}{@{}>{\leavevmode\bfseries\color{navy}}p{27mm}>{\leavevmode}X@{}}
  Tested against & #1\\
  Missing in Julia & #2\\
  Package & #3\\
  Result & #4
  \end{tabularx}
  \end{tcolorbox}}

% Boxes
\newtcolorbox{pbox}[2][]{enhanced, colback=pale, colframe=amber, boxrule=0pt,
  leftrule=3pt, sharp corners, left=7pt, right=7pt, top=6pt, bottom=6pt,
  fontupper=\small\raggedright, before upper={{\color{navy}\normalsize\bfseries #2}\par\smallskip}, #1}
\newtcolorbox{nbox}[2][]{enhanced, colback=navy, colframe=navy, boxrule=0pt,
  sharp corners, left=8pt, right=8pt, top=7pt, bottom=7pt, coltext=white,
  fontupper=\small\raggedright, before upper={{\color{amber}\normalsize\bfseries #2}\par\smallskip}, #1}

% Value tables: published beside Julia
\newenvironment{valuetable}[1]{%
  \par\small\tabularx{\linewidth}{#1}\toprule}{\bottomrule\endtabularx\par}
\newcommand{\ok}{{\color{agree}\faCheck}}
\newcommand{\no}{{\color{differ}\faTimes}}

% End of the document: credit and licence
\newcommand{\nmcredit}{%
  \par\vfill
  \begin{minipage}[b]{0.78\linewidth}\footnotesize\color{soft}
  A Pollis Toolbox Guide, written from the panels of the toolbox.\par
  \textcopyright\ \vignetteauthor\enspace·\enspace Licensed under \vignettelicense\enspace·\enspace\vignettedate
  \end{minipage}\hfill
  \begin{minipage}[b]{0.18\linewidth}\raggedleft\pollislogo{9mm}\end{minipage}}
.
\documentclass[10pt,a4paper]{article}

\usepackage[a4paper,left=18mm,right=18mm,top=15mm,bottom=17mm,footskip=8mm]{geometry}
\usepackage{iftex}
\ifPDFTeX\usepackage[T1]{fontenc}\fi
\usepackage[default,semibold]{sourcesanspro}
\usepackage{textcomp}
\usepackage[italic]{mathastext}
\usepackage{amsmath}
\usepackage{microtype}
\usepackage[table]{xcolor}
\usepackage{tikz}
\usepackage[most]{tcolorbox}
\usepackage{fontawesome5}
\usepackage{tabularx,booktabs}
\usepackage{enumitem}
\usepackage{fancyhdr}
\usepackage[hidelinks]{hyperref}

% ---------------------------------------------------------------- colours (as in the brochure)
\definecolor{navy}{HTML}{1C3A5E}
\definecolor{pblue}{HTML}{3B78B5}
\definecolor{amber}{HTML}{F2B33D}
\definecolor{deepamber}{HTML}{D9902A}
\definecolor{pale}{HTML}{EEF3F9}
\definecolor{ink}{HTML}{2A2F36}
\definecolor{logonavy}{HTML}{173D63}
\definecolor{logodark}{HTML}{C98A1E}
\definecolor{logoamber}{HTML}{F2B446}
\definecolor{logocream}{HTML}{FEEFD2}
\definecolor{soft}{HTML}{5E6773}
\definecolor{agree}{HTML}{2F7D4F}
\definecolor{differ}{HTML}{B23B3B}
\color{ink}

\setlength{\parindent}{0pt}
\setlength{\parskip}{4pt}
\setlist{nosep,leftmargin=*,itemsep=2pt}
\renewcommand{\arraystretch}{1.15}

\pagestyle{fancy}
\fancyhf{}
\renewcommand{\headrulewidth}{0pt}
\fancyfoot[L]{\footnotesize\color{soft}\textbf{Pollis}\enspace Toolbox Guide\enspace·\enspace\nmtheme}
\fancyfoot[R]{\footnotesize\color{soft}\thepage}
\newcommand{\nmtheme}{}
\newcommand{\vignetteauthor}{}
\newcommand{\vignettedate}{\today}
\newcommand{\vignettelicense}{\href{https://creativecommons.org/licenses/by-sa/4.0/}{CC BY-SA 4.0}}

% The Pollis logo, redrawn so it prints sharply at any size
\newcommand{\pollislogo}[1]{%
  \begin{tikzpicture}[x=#1/512,y=-#1/512,baseline=(current bounding box.center)]
    \fill[logonavy] (0,0) rectangle (512,512);
    \fill[logodark] (36,422) -- (104,256) -- (172,422) -- cycle;
    \fill[logoamber] (104,422) -- (203,180) -- (302,422) -- cycle;
    \fill[logocream] (203,422) -- (339,89) -- (476,422) -- cycle;
  \end{tikzpicture}}

% Top of the first page: logo, name; then the kicker, title and a one-line promise
\newcommand{\nmhead}[3]{%
  \renewcommand{\nmtheme}{#1}%
  \thispagestyle{fancy}%
  \begin{minipage}[c]{13mm}\pollislogo{11mm}\end{minipage}%
  \begin{minipage}[c]{\dimexpr\linewidth-13mm\relax}
    {\color{navy}\Large\bfseries Pollis}\enspace{\color{soft}Toolbox Guide}\hfill
    {\color{deepamber}\small\bfseries\MakeUppercase{#1}}
  \end{minipage}\par\vspace{7mm}
  {\color{navy}\fontsize{22}{26}\selectfont\bfseries #2}\par\vspace{4pt}
  {\color{soft}\large #3}\par\vspace{5mm}}

\newcommand{\nmsection}[1]{\par\vspace{5pt}{\color{navy}\large\bfseries #1}\par\vspace{1pt}}
\newcommand{\nmsub}[1]{\par\vspace{2pt}{\color{deepamber}\bfseries #1}\par}
\newcommand{\gt}{\,{\color{deepamber}\bfseries ›}\,}
\newcommand{\menupath}[1]{{\bfseries\color{navy}#1}}
\newcommand{\fn}[1]{\texttt{\small #1}}
% Citations: \nmcite{key}{Breiman (2001)} links to the entry \nmref{key}{...} in the References
\newcommand{\nmcite}[2]{\hyperlink{ref:#1}{\color{pblue}#2}}
\newcommand{\nmref}[2]{\raisebox{1.1em}[0pt][0pt]{\hypertarget{ref:#1}{}}#2\par}
\newcommand{\nmdoi}[1]{\href{https://doi.org/#1}{\color{pblue}doi:#1}}
% Where the package is: its name and a link (its GitHub repository, once published)
\newcommand{\package}[2]{\href{#2}{\fn{#1.jl}}}

% Key facts at the top: source, missing in Julia, package, result
\newcommand{\factsheet}[4]{%
  \begin{tcolorbox}[enhanced, colback=pale, colframe=pale, boxrule=0pt, sharp corners,
    left=7pt, right=7pt, top=5pt, bottom=5pt, fontupper=\small]
  \begin{tabularx}{\linewidth}{@{}>{\leavevmode\bfseries\color{navy}}p{27mm}>{\leavevmode}X@{}}
  Tested against & #1\\
  Missing in Julia & #2\\
  Package & #3\\
  Result & #4
  \end{tabularx}
  \end{tcolorbox}}

% Boxes
\newtcolorbox{pbox}[2][]{enhanced, colback=pale, colframe=amber, boxrule=0pt,
  leftrule=3pt, sharp corners, left=7pt, right=7pt, top=6pt, bottom=6pt,
  fontupper=\small\raggedright, before upper={{\color{navy}\normalsize\bfseries #2}\par\smallskip}, #1}
\newtcolorbox{nbox}[2][]{enhanced, colback=navy, colframe=navy, boxrule=0pt,
  sharp corners, left=8pt, right=8pt, top=7pt, bottom=7pt, coltext=white,
  fontupper=\small\raggedright, before upper={{\color{amber}\normalsize\bfseries #2}\par\smallskip}, #1}

% Value tables: published beside Julia
\newenvironment{valuetable}[1]{%
  \par\small\tabularx{\linewidth}{#1}\toprule}{\bottomrule\endtabularx\par}
\newcommand{\ok}{{\color{agree}\faCheck}}
\newcommand{\no}{{\color{differ}\faTimes}}

% End of the document: credit and licence
\newcommand{\nmcredit}{%
  \par\vfill
  \begin{minipage}[b]{0.78\linewidth}\footnotesize\color{soft}
  A Pollis Toolbox Guide, written from the panels of the toolbox.\par
  \textcopyright\ \vignetteauthor\enspace·\enspace Licensed under \vignettelicense\enspace·\enspace\vignettedate
  \end{minipage}\hfill
  \begin{minipage}[b]{0.18\linewidth}\raggedleft\pollislogo{9mm}\end{minipage}}
.
\documentclass[10pt,a4paper]{article}

\usepackage[a4paper,left=18mm,right=18mm,top=15mm,bottom=17mm,footskip=8mm]{geometry}
\usepackage{iftex}
\ifPDFTeX\usepackage[T1]{fontenc}\fi
\usepackage[default,semibold]{sourcesanspro}
\usepackage{textcomp}
\usepackage[italic]{mathastext}
\usepackage{amsmath}
\usepackage{microtype}
\usepackage[table]{xcolor}
\usepackage{tikz}
\usepackage[most]{tcolorbox}
\usepackage{fontawesome5}
\usepackage{tabularx,booktabs}
\usepackage{enumitem}
\usepackage{fancyhdr}
\usepackage[hidelinks]{hyperref}

% ---------------------------------------------------------------- colours (as in the brochure)
\definecolor{navy}{HTML}{1C3A5E}
\definecolor{pblue}{HTML}{3B78B5}
\definecolor{amber}{HTML}{F2B33D}
\definecolor{deepamber}{HTML}{D9902A}
\definecolor{pale}{HTML}{EEF3F9}
\definecolor{ink}{HTML}{2A2F36}
\definecolor{logonavy}{HTML}{173D63}
\definecolor{logodark}{HTML}{C98A1E}
\definecolor{logoamber}{HTML}{F2B446}
\definecolor{logocream}{HTML}{FEEFD2}
\definecolor{soft}{HTML}{5E6773}
\definecolor{agree}{HTML}{2F7D4F}
\definecolor{differ}{HTML}{B23B3B}
\color{ink}

\setlength{\parindent}{0pt}
\setlength{\parskip}{4pt}
\setlist{nosep,leftmargin=*,itemsep=2pt}
\renewcommand{\arraystretch}{1.15}

\pagestyle{fancy}
\fancyhf{}
\renewcommand{\headrulewidth}{0pt}
\fancyfoot[L]{\footnotesize\color{soft}\textbf{Pollis}\enspace Toolbox Guide\enspace·\enspace\nmtheme}
\fancyfoot[R]{\footnotesize\color{soft}\thepage}
\newcommand{\nmtheme}{}
\newcommand{\vignetteauthor}{}
\newcommand{\vignettedate}{\today}
\newcommand{\vignettelicense}{\href{https://creativecommons.org/licenses/by-sa/4.0/}{CC BY-SA 4.0}}

% The Pollis logo, redrawn so it prints sharply at any size
\newcommand{\pollislogo}[1]{%
  \begin{tikzpicture}[x=#1/512,y=-#1/512,baseline=(current bounding box.center)]
    \fill[logonavy] (0,0) rectangle (512,512);
    \fill[logodark] (36,422) -- (104,256) -- (172,422) -- cycle;
    \fill[logoamber] (104,422) -- (203,180) -- (302,422) -- cycle;
    \fill[logocream] (203,422) -- (339,89) -- (476,422) -- cycle;
  \end{tikzpicture}}

% Top of the first page: logo, name; then the kicker, title and a one-line promise
\newcommand{\nmhead}[3]{%
  \renewcommand{\nmtheme}{#1}%
  \thispagestyle{fancy}%
  \begin{minipage}[c]{13mm}\pollislogo{11mm}\end{minipage}%
  \begin{minipage}[c]{\dimexpr\linewidth-13mm\relax}
    {\color{navy}\Large\bfseries Pollis}\enspace{\color{soft}Toolbox Guide}\hfill
    {\color{deepamber}\small\bfseries\MakeUppercase{#1}}
  \end{minipage}\par\vspace{7mm}
  {\color{navy}\fontsize{22}{26}\selectfont\bfseries #2}\par\vspace{4pt}
  {\color{soft}\large #3}\par\vspace{5mm}}

\newcommand{\nmsection}[1]{\par\vspace{5pt}{\color{navy}\large\bfseries #1}\par\vspace{1pt}}
\newcommand{\nmsub}[1]{\par\vspace{2pt}{\color{deepamber}\bfseries #1}\par}
\newcommand{\gt}{\,{\color{deepamber}\bfseries ›}\,}
\newcommand{\menupath}[1]{{\bfseries\color{navy}#1}}
\newcommand{\fn}[1]{\texttt{\small #1}}
% Citations: \nmcite{key}{Breiman (2001)} links to the entry \nmref{key}{...} in the References
\newcommand{\nmcite}[2]{\hyperlink{ref:#1}{\color{pblue}#2}}
\newcommand{\nmref}[2]{\raisebox{1.1em}[0pt][0pt]{\hypertarget{ref:#1}{}}#2\par}
\newcommand{\nmdoi}[1]{\href{https://doi.org/#1}{\color{pblue}doi:#1}}
% Where the package is: its name and a link (its GitHub repository, once published)
\newcommand{\package}[2]{\href{#2}{\fn{#1.jl}}}

% Key facts at the top: source, missing in Julia, package, result
\newcommand{\factsheet}[4]{%
  \begin{tcolorbox}[enhanced, colback=pale, colframe=pale, boxrule=0pt, sharp corners,
    left=7pt, right=7pt, top=5pt, bottom=5pt, fontupper=\small]
  \begin{tabularx}{\linewidth}{@{}>{\leavevmode\bfseries\color{navy}}p{27mm}>{\leavevmode}X@{}}
  Tested against & #1\\
  Missing in Julia & #2\\
  Package & #3\\
  Result & #4
  \end{tabularx}
  \end{tcolorbox}}

% Boxes
\newtcolorbox{pbox}[2][]{enhanced, colback=pale, colframe=amber, boxrule=0pt,
  leftrule=3pt, sharp corners, left=7pt, right=7pt, top=6pt, bottom=6pt,
  fontupper=\small\raggedright, before upper={{\color{navy}\normalsize\bfseries #2}\par\smallskip}, #1}
\newtcolorbox{nbox}[2][]{enhanced, colback=navy, colframe=navy, boxrule=0pt,
  sharp corners, left=8pt, right=8pt, top=7pt, bottom=7pt, coltext=white,
  fontupper=\small\raggedright, before upper={{\color{amber}\normalsize\bfseries #2}\par\smallskip}, #1}

% Value tables: published beside Julia
\newenvironment{valuetable}[1]{%
  \par\small\tabularx{\linewidth}{#1}\toprule}{\bottomrule\endtabularx\par}
\newcommand{\ok}{{\color{agree}\faCheck}}
\newcommand{\no}{{\color{differ}\faTimes}}

% End of the document: credit and licence
\newcommand{\nmcredit}{%
  \par\vfill
  \begin{minipage}[b]{0.78\linewidth}\footnotesize\color{soft}
  A Pollis Toolbox Guide, written from the panels of the toolbox.\par
  \textcopyright\ \vignetteauthor\enspace·\enspace Licensed under \vignettelicense\enspace·\enspace\vignettedate
  \end{minipage}\hfill
  \begin{minipage}[b]{0.18\linewidth}\raggedleft\pollislogo{9mm}\end{minipage}}

\renewcommand{\vignetteauthor}{Pollis}
\renewcommand{\vignettedate}{4 November 2026}
\renewcommand{\vignettelicense}{AGPL-3.0}

% The state of an entry of the toolbox
\newcommand{\ready}{\leavevmode{\color{agree}\bfseries\faCheck\enspace Ready}}
\newcommand{\soon}{\leavevmode{\color{deepamber}\bfseries\faHourglassHalf\enspace In preparation}}

\begin{document}

\nmhead{Self-Study: Learn and Practise}
  {Self-Study Guide}
  {Turn your slides, notes and PDFs into what helps you study.}

\begin{nbox}{Self-Study: Learn and Practise}
Turn your slides, notes and PDFs into what helps you study: more problems with their solutions, a topic
explained in depth, a summary of a paper, prose and audio from slides, or slides from a long text. For
students, and for teachers preparing their classes.
\end{nbox}

Most toolboxes bring a field to Pollis: finance, economics, mathematics. This one brings a way of working.
It eases the chores of studying, such as finding more exercises, making sense of a dense paragraph or
getting through a long paper, so that more of your time goes to learning itself. Pollis AI+ is its study companion:
it writes the problems, the explanations, the summaries and the slides from the material you give it, and
Julia computes every number that has to be right.

\nmsection{What the Toolbox Does}
{\small
\begin{tabularx}{\linewidth}{@{}>{\raggedright\arraybackslash}p{44mm}X>{\raggedright\arraybackslash}p{27mm}@{}}
\toprule
\textbf{Entry} & \textbf{What you get} & \textbf{State} \\\midrule
Practice & Multiple-choice exercises on descriptive statistics, hypothesis tests and linear regression, with new data for each seed & \ready \\
Write an Exercise & Your own exercises, computed from data or conceptual, with a version for each seed & \ready \\
Make More Problems & More problems like those of your slides or notes, with their solutions & \soon \\
Explain a Topic in Depth & A paragraph or a topic expanded, with each concept detailed & \soon \\
Summarise a Paper & The summary of a paper in PDF, as a PDF & \soon \\
Slides into Prose and Audio & Slides turned into continuous text, and an audio version of it & \soon \\
Text into Slides & A long text turned into a few slides & \soon \\\bottomrule
\end{tabularx}}

\nmsection{Install the Toolbox}
\begin{enumerate}[label=\arabic*.]
  \item Open the Extensions pane, in the activity bar on the left. The \menupath{Toolboxes} section, at the
    top, lists the toolboxes you have not installed.
  \item Choose \menupath{Install} beside \menupath{Self-Study: Learn and Practise}, then \menupath{Reload Now}
    when Pollis asks.
  \item On the \menupath{Toolbox Panels} page that opens, tick the panels you want and choose
    \menupath{Publish}. They appear in \menupath{Extend \gt Self-Study: Learn and Practise}.
\end{enumerate}
This guide and the toolbox's vignette are in \menupath{Help \gt Documentation \gt Toolboxes}, and Ask Pollis
answers questions about the toolbox from them.

\nmsection{Practise a Topic}
\menupath{Extend \gt Self-Study: Learn and Practise \gt Practice} has a tab for each topic:
\menupath{Descriptive Statistics}, \menupath{Hypothesis Tests} and \menupath{Linear Regression}. Each
exercise asks one question with four choices. Its data come from a random number generator started from
the \menupath{Seed}, and its answers are computed from those data, so they are always right. The wrong choices
are the usual mistakes, such as dividing by $n$ instead of $n-1$, or a one-sided p-value where a two-sided
one is asked for.
\begin{enumerate}[label=\arabic*.]
  \item Choose a tab and an \menupath{Exercise}.
  \item Run the Example Code: the question and the choices A to D appear in the Julia REPL, and the data stay
    there, so you can work the answer out in Julia as well as on paper.
  \item Pick your answer under \menupath{Your Answer}.
  \item Under Next Steps, \menupath{Check My Answer} says whether it is correct and explains it;
    \menupath{Show the Answer} explains it without your answer.
  \item Change the \menupath{Seed} and run again: other data, another order of the choices, the same skill.
\end{enumerate}

\nmsection{Write Your Own Exercises}
\menupath{Write an Exercise} turns an exercise of your own into one that Julia computes. An exercise is a Julia
function: from a random number generator it returns the question, the choices with the correct one first,
and the explanation.
\begin{itemize}
  \item \menupath{Computed from Data} draws the numbers from the seed and computes the choices from them: for a
    statistic, a p-value or a prediction.
  \item \menupath{Conceptual} keeps the text fixed and shuffles the choices: for choosing a method or reading a
    result.
\end{itemize}
Give the exercise a \menupath{Title}, a \menupath{First Seed} and the number of \menupath{Versions} to preview.
Before you use it, three Next Steps check it: \menupath{The Correct Answer Across Versions} shows the range of
the answer, \menupath{Choices That Coincide} looks for versions where two choices come out the same, and
\menupath{Try a Version} shows one as a student sees it. The exercises you write are kept in
\fn{\textasciitilde/.pollis/practice}, and Practice offers them beside its own.

\nmsection{Pollis AI+ as Your Study Companion}
The five entries below take material you already have and give back material that helps you study. You
choose the file, Pollis AI+ does the work, and the result is a file of your own, to read, print or edit.
\begin{pbox}{What you give, what you get}
\begin{tabularx}{\linewidth}{@{}>{\raggedright\arraybackslash\bfseries}p{36mm}>{\raggedright\arraybackslash}p{46mm}X@{}}
Make More Problems & Slides, notes or a PDF with a few problems, and how many more you want & New problems on the same topic and at the same level, each with its full solution, as a PDF; those computed from data can also become exercises of Write an Exercise \\[3pt]
Explain a Topic in Depth & A paragraph, an area of a PDF page, or the name of a topic & The concepts of the paragraph detailed one by one, with the steps of each argument and an example, as a PDF \\[3pt]
Summarise a Paper & A paper in PDF & Its question, data, method, results and limits, with the pages they come from, as a PDF \\[3pt]
Slides into Prose and Audio & A set of slides, in PDF or LaTeX & The slides as continuous text, as a PDF, and an audio file that reads it \\[3pt]
Text into Slides & A long text: a PDF, a LaTeX or Markdown file, or a notebook & A few slides with its main points, as a PDF, with their LaTeX source to edit
\end{tabularx}
\end{pbox}
A summary or an explanation is only as good as its source: read it beside the original, and check the
numbers that matter. The PDF Viewer of Pollis opens both side by side, and its \menupath{Select Area} takes
the text, a table or a picture of any part of a page.

\nmsection{For Teachers}
The same entries prepare a class. Write an Exercise gives each seed its own version, with other data and
another order of the choices, so every student, or every attempt, can get a different one while practising
the same skill. Make More Problems extends the exercises of a set of slides; Text into Slides and Slides
into Prose and Audio give a lecture in the form each student prefers to study from.

\nmsection{Coming Next}
The entries in preparation come in this order, each added to the menu of the toolbox as it is ready:
\begin{enumerate}[label=\arabic*.]
  \item \textbf{Make More Problems}, which builds on Write an Exercise.
  \item \textbf{Summarise a Paper}, which builds on the PDF Viewer.
  \item \textbf{Explain a Topic in Depth}, from a paragraph or an area of a page.
  \item \textbf{Text into Slides}, compiled with PollisTeX, the \LaTeX\ of Pollis.
  \item \textbf{Slides into Prose and Audio}, with the audio read by a voice model.
\end{enumerate}
Each one is described in this guide as it comes, and the toolbox's version tells you which you have.

\nmcredit
\end{document}
